\chapter{Coordinate Systems on the Celestial Sphere}

Chapter~1 established the temporal frameworks -- solar time, sidereal time, and precession -- against which every astronomical observation must be referenced. To specify the locations of stars, planets, and galaxies, astronomers project the three-dimensional Universe onto the two-dimensional surface of an imaginary sphere of arbitrary radius centered on the observer: the \emph{celestial sphere}. Because astronomical measurements are performed from different vantage points (from an observatory on Earth, from the moving Earth relative to the Sun, or across the Milky Way), several coordinate systems are employed. In this chapter, we derive the four principal celestial coordinate systems --- Equatorial, Horizontal, Ecliptic, and Galactic --- and establish the rigorous vector transformations relating them.

\section{The Equatorial Coordinate System}

The \emph{equatorial coordinate system} extends the Earth's terrestrial latitude-longitude grid onto the celestial sphere: the celestial equator is the projection of the Earth's equator, the north and south celestial poles ($\mathrm{NCP}, \mathrm{SCP}$) are the projections of the Earth's rotation axis, and coordinates are fixed relative to the distant stars.

\begin{definition}[Right Ascension and Declination]
	\emph{Declination} $\delta \in [-90^\circ, +90^\circ]$ is the angular distance north ($+$) or south ($-$) of the celestial equator. \emph{Right ascension} $\alpha \in [0^{\rm h}, 24^{\rm h})$ ($0^\circ \le \alpha < 360^\circ$) is the angle measured eastward along the celestial equator from the vernal equinox $\hat\gamma$.
\end{definition}

\begin{figure}[htbp]
	\centering
	\includegraphics[width=0.55\linewidth]{declination_celestial_coordinates_en_bviPcZn.pdf}
	\caption{The equatorial Cartesian frame: the $z$-axis lies along the Earth's rotation axis ($\mathrm{NCP}$), the $xy$-plane is the celestial equator, and the $x$-axis points toward the vernal equinox $\hat\gamma$.}
	\label{fig:declination}
\end{figure}

In the equatorial Cartesian frame centered on the Earth, a source with coordinates $(\alpha, \delta)$ has the unit position vector:
\begin{keyresult}
	\begin{equation}
		\hat{\mathbf r}_{\rm eq} = \begin{pmatrix} \cos\delta\cos\alpha \\ \cos\delta\sin\alpha \\ \sin\delta \end{pmatrix}.
		\label{eq:unit_vector_eq}
	\end{equation}
\end{keyresult}

Because the equatorial grid is attached to the celestial sphere, stellar catalog coordinates $(\alpha, \delta)$ remain nearly constant on human timescales (subject only to proper motion and precession). The apparent daily rising and setting of celestial objects reflects the rotation of the observer's horizon underneath this fixed background.

\subsection{Coordinate Epochs}
\label{sec:coordinate_epochs}

Because the equinox $\hat\gamma$ that anchors the $(\alpha,\delta)$ grid precesses at the rate derived in Chapter~1 (Eq.~\eqref{eq:precession_rate}), a catalogued position $(\alpha,\delta)$ is meaningless without a stated \emph{epoch} -- the instant at which the equatorial grid used to define it is evaluated. The current standard is J2000.0; older catalogs commonly used B1950.0. Some source names encode the epoch directly in the identifier itself: for example, the pulsar $\mathrm{PSR\ J1939{+}2134}$ is named for its J2000 coordinates $\alpha \approx 19^{\rm h}39^{\rm m}$, $\delta \approx +21^\circ34'$.

Converting a cataloged position from one epoch to another in general requires \emph{two} independent corrections: the purely geometric precession of the reference grid (Eqs.~\eqref{eq:precess_alpha}--\eqref{eq:precess_delta}), and the star's own physical proper motion $(\mu_\alpha^*, \mu_\delta)$ across the sky (derived in Chapter~3). The two effects are independent and additive over the (typically short, $\Delta t \lesssim$ decades) baseline of a proper-motion correction, but only precession is required to transform a fixed grid definition (e.g.\ B1950.0 $\to$ J2000.0) when the object itself is treated as stationary.

\section{The Horizontal (Alt-Az) Coordinate System and the Hour Angle}

The \emph{horizontal} coordinate system is attached to the observer's local horizon and zenith, co-rotating with the Earth.

\begin{definition}[Altitude, Azimuth, and Hour Angle]
	\emph{Altitude} $h \in [-90^\circ, +90^\circ]$ is the angle of a celestial object above the local horizon (zenith at $h = 90^\circ$, nadir at $h = -90^\circ$). \emph{Azimuth} $A \in [0^\circ, 360^\circ)$ is the angle along the horizon measured from North ($0^\circ$) through East ($90^\circ$), South ($180^\circ$), and West ($270^\circ$). The \emph{Hour Angle} $H$ is the angular distance measured along the celestial equator from the local meridian to the object's hour circle:
	\begin{equation}
		H \equiv \LST - \alpha,
		\label{eq:hour_angle_formula}
	\end{equation}
	with $H > 0$ after meridian transit (moving westward).
\end{definition}

\begin{figure}[htbp]
	\centering
	\includegraphics[width=0.75\linewidth]{altitude_azimuth_en.pdf}
	\caption{The horizontal coordinate system: the local meridian is the great circle passing through the North cardinal point, the Zenith, and the South cardinal point.}
	\label{fig:altaz}
\end{figure}

\section{Derivation of the Horizontal $\leftrightarrow$ Equatorial Transformations}

We now derive the transformations relating horizontal coordinates $(h, A)$ and equatorial coordinates $(\delta, H)$ for an observer at terrestrial latitude $\phi_0$. We express the unit vector toward a celestial source in two orthonormal coordinate frames sharing the same origin and relate them via an orthogonal rotation matrix.

\subsection{Source Vector in the Meridian-Aligned Equatorial Frame}

Align the Cartesian $x$-axis with the observer's \emph{local celestial meridian} at the instant of observation ($H = 0$). In this frame, the source's equatorial coordinate along the equator is $-H$:
\begin{equation}
	\hat{\mathbf s}_{\rm eq} = \begin{pmatrix} \cos\delta\cos(-H) \\ \cos\delta\sin(-H) \\ \sin\delta \end{pmatrix} = \begin{pmatrix} \cos\delta\cos H \\ -\cos\delta\sin H \\ \sin\delta \end{pmatrix}.
	\label{eq:unit_meridian_frame}
\end{equation}

\subsection{The Local Horizontal Basis}

The cardinal horizontal directions --- North ($\hat{\mathbf n}$), East ($\hat{\mathbf e}$), and Zenith ($\hat{\mathbf z}$) --- expressed in the meridian equatorial frame for latitude $\phi_0$ are:
\begin{itemize}
	\item \textbf{East ($\hat{\mathbf e}$):} Along the celestial equator, $90^\circ$ east of meridian: $(\delta, H) = (0, -90^\circ)$.
	\item \textbf{Zenith ($\hat{\mathbf z}$):} Overhead on the meridian: $(\delta, H) = (\phi_0, 0^\circ)$.
	\item \textbf{North ($\hat{\mathbf n}$):} Point on the horizon directly North, $90^\circ$ from zenith along the meridian: $(\delta, H) = (90^\circ + \phi_0, 0^\circ)$.
\end{itemize}

\begin{figure}[htbp]
	\begin{minipage}[b]{0.48\textwidth}
		\centering
		Substituting into Eq.~\eqref{eq:unit_meridian_frame}:
		\begin{align}
			\hat{\mathbf n} &= \begin{pmatrix} -\sin\phi_0 \\ 0 \\ \cos\phi_0 \end{pmatrix}, \label{eq:basis_n} \\
			\hat{\mathbf e} &= \begin{pmatrix} 0 \\ 1 \\ 0 \end{pmatrix}, \label{eq:basis_e} \\
			\hat{\mathbf z} &= \begin{pmatrix} \cos\phi_0 \\ 0 \\ \sin\phi_0 \end{pmatrix}. \label{eq:basis_z}
		\end{align}
	\end{minipage}
	\hfill
	\begin{minipage}[b]{0.48\textwidth}
		\centering
		\includegraphics[width=0.7\linewidth]{local north and east.png}
		\caption{Local horizontal basis directions $\{\hat{\mathbf n}, \hat{\mathbf e}, \hat{\mathbf z}\}$.}
		\label{fig:local_coordinates}
	\end{minipage}
\end{figure}

\subsection{Transformation Equations}

Any unit vector can be decomposed in the horizontal basis $(h, A)$ as:
\begin{equation}
	\hat{\mathbf s} = (\cos h\cos A)\,\hat{\mathbf n} + (\cos h\sin A)\,\hat{\mathbf e} + (\sin h)\,\hat{\mathbf z}.
\end{equation}
Using the orthogonal rotation matrix $\mathbf R$:
\begin{equation}
	\begin{pmatrix} \cos h\cos A \\ \cos h\sin A \\ \sin h \end{pmatrix} = \begin{pmatrix} -\sin\phi_0 & 0 & \cos\phi_0 \\ 0 & 1 & 0 \\ \cos\phi_0 & 0 & \sin\phi_0 \end{pmatrix} \begin{pmatrix} \cos\delta\cos H \\ -\cos\delta\sin H \\ \sin\delta \end{pmatrix}.
	\label{eq:eq_to_hor_matrix}
\end{equation}
Multiplying out row by row yields the fundamental transformation relations:
\begin{keyresult}
	\begin{align}
		\sin h &= \sin\phi_0\sin\delta + \cos\phi_0\cos\delta\cos H, \label{eq:sinh_master} \\
		\cos h\sin A &= -\cos\delta\sin H, \label{eq:cosh_sinA} \\
		\cos h\cos A &= -\sin\phi_0\cos\delta\cos H + \cos\phi_0\sin\delta. \label{eq:cosh_cosA}
	\end{align}
\end{keyresult}
Dividing Eq.~\eqref{eq:cosh_sinA} by Eq.~\eqref{eq:cosh_cosA} gives the azimuth:
\begin{equation}
	\tan A = \frac{-\cos\delta\sin H}{-\sin\phi_0\cos\delta\cos H + \cos\phi_0\sin\delta} = \frac{\sin H}{\sin\phi_0\cos H - \cos\phi_0\tan\delta}.
	\label{eq:tanA_master}
\end{equation}

\begin{example}[Sirius at LST $22^{\rm h}00^{\rm m}$ from IISER TVM]
	For Sirius ($\alpha = 06^{\rm h}45^{\rm m}09^{\rm s}$, $\delta = -16.716^\circ$) observed from IISER TVM ($\phi_0 = 8.6833^\circ\text{N}$) at $\LST = 22^{\rm h}00^{\rm m}00^{\rm s}$:
	\[
	H = \LST - \alpha = 15^{\rm h}14^{\rm m}51^{\rm s} \approx 228.71^\circ.
	\]
	From Eq.~\eqref{eq:sinh_master}:
	\[
	\sin h = \sin(8.6833^\circ)\sin(-16.716^\circ) + \cos(8.6833^\circ)\cos(-16.716^\circ)\cos(228.71^\circ) = -0.668 \implies h \approx -41.9^\circ.
	\]
	Sirius is well below the horizon. From Eq.~\eqref{eq:tanA_master}, taking into account signs of numerator and denominator, $A \approx 104.74^\circ$.
\end{example}

\section{Rising, Setting, and Circumpolar Stars}

A celestial object rises or sets when its geometric center crosses the true horizon ($h = 0^\circ$). Setting $h = 0$ in Eq.~\eqref{eq:sinh_master}:
\begin{keyresult}
	\begin{equation}
		\cos H_0 = -\tan\phi_0\tan\delta.
		\label{eq:rise_set_master}
	\end{equation}
\end{keyresult}
This yields two solutions:
\begin{itemize}
	\item \textbf{Setting:} $H = +H_0 = \arccos(-\tan\phi_0\tan\delta) > 0$.
	\item \textbf{Rising:} $H = -H_0 = -\arccos(-\tan\phi_0\tan\delta) < 0$.
\end{itemize}
The local sidereal times of rising and setting are $\LST_{\rm rise} = \alpha - H_0$ and $\LST_{\rm set} = \alpha + H_0$.

\begin{figure}[htbp]
	\centering
	\includegraphics[width=0.55\linewidth]{rising-setting.png}
	\caption{Diurnal motion of stars along small circles of constant declination.}
	\label{fig:rising_setting}
\end{figure}

\paragraph{Circumpolar and Never-Rising Conditions.}
The quantity $|\tan\phi_0\tan\delta|$ determines visibility:
\begin{itemize}
	\item If $|\tan\phi_0\tan\delta| \le 1$: The star rises and sets daily.
	\item If $\tan\phi_0\tan\delta > 1$ ($\delta > 90^\circ - |\phi_0|$ in Northern hemisphere): The object is \emph{circumpolar} and never sets ($\cos H_0 < -1$, $h > 0$ always).
	\item If $\tan\phi_0\tan\delta < -1$ ($\delta < -(90^\circ - |\phi_0|)$): The object \emph{never rises} above the horizon.
\end{itemize}

\section{Zenith Angle at Meridian Transit and Local Solar Noon}

\begin{definition}[Zenith Angle]
	The \emph{zenith angle} $z$ is the angular distance from the observer's zenith:
	\begin{equation}
		z = 90^\circ - h.
		\label{eq:zenith_angle_def}
	\end{equation}
\end{definition}

At upper meridian transit (culmination), $H = 0$, so $\cos H = 1$. Eq.~\eqref{eq:sinh_master} simplifies to:
\[
\sin h_{\rm max} = \sin\phi_0\sin\delta + \cos\phi_0\cos\delta = \cos(\phi_0 - \delta).
\]
Since $\sin h = \cos z$, this yields the fundamental relation:
\begin{keyresult}
	\begin{equation}
		z_{\rm min} = |\phi_0 - \delta|.
		\label{eq:zenith_angle_noon}
	\end{equation}
\end{keyresult}
At local solar noon ($\delta = \delta_\odot$), the Sun's zenith angle is simply $z_{\rm noon} = |\phi_0 - \delta_\odot|$.

\begin{example}[Summer Solstice Noon at IISER TVM]
	At the June solstice, $\delta_\odot = +23.44^\circ$. For IISER TVM ($\phi_0 = 8.6833^\circ\text{N}$):
	\[
	z_{\rm noon} = |8.6833^\circ - 23.44^\circ| = 14.76^\circ.
	\]
	The Sun culminates $14.76^\circ$ north of zenith.
\end{example}

\section{Atmospheric Refraction and Stellar Aberration}

\subsection{Atmospheric Refraction}
As starlight enters the Earth's atmosphere, the increasing refractive index bends the ray toward the zenith, making objects appear \emph{higher} than their true geometric altitude:
\begin{equation}
	h_{\rm apparent} = h_{\rm true} + R(h).
\end{equation}
To see where the numerical coefficient in this bending comes from, again use the plane-parallel approximation: treat the atmosphere as a single slab of refractive index $n$ (vacuum, $n=1$, above it). Snell's law relates the true zenith angle $z$ (outside the atmosphere) to the apparent zenith angle $z' = z - R$ seen by the observer at the surface:
\begin{equation}
	\sin z = n\sin z' = n\sin(z - R).
\end{equation}
Expanding $\sin(z-R) \approx \sin z - R\cos z$ for small $R$:
\begin{equation}
	\sin z \approx n\sin z - nR\cos z \quad\Longrightarrow\quad R \approx \frac{(n-1)\sin z}{n\cos z} \approx (n-1)\tan z,
	\label{eq:refraction_derivation}
\end{equation}
using $n\approx 1$ in the denominator. With the refractive index of air at the surface $n - 1 \approx 2.8\times10^{-4}$ (a measured property of air, not something derivable from geometry alone), and converting to arcseconds ($1\,\text{rad} = 206{,}265''$):
\begin{keyresult}
	\begin{equation}
		R \approx (n-1)\tan z \times 206{,}265'' \approx 58.2''\tan z = 58.2''\cot h_{\rm true} \qquad (h_{\rm true} > 15^\circ).
		\label{eq:refraction_cot}
	\end{equation}
\end{keyresult}
Because $n-1$ itself depends on air density (hence on local temperature and pressure), this coefficient is only a standard-atmosphere reference value; the empirical formula below builds in an explicit correction for the actual $(P,T)$ at the telescope. Near the horizon, the single-slab approximation breaks down entirely (as $z\to90^\circ$ the ray grazes the atmosphere over a curved path), so Bennett's formula provides sub-arcsecond accuracy there instead, as an empirical fit rather than a derived result:
\begin{equation}
	R[\text{arcmin}] \approx \frac{1.02}{\tan\!\left(h + \dfrac{10.3^\circ}{h + 5.11^\circ}\right)}\times\left(\frac{P}{1010\,\text{hPa}}\right)\left(\frac{283}{273 + T[^\circ\text{C}]}\right).
	\label{eq:bennett_formula}
\end{equation}
At the horizon ($h = 0^\circ$), $R \approx 34'$ (greater than the Sun's angular diameter $\sim 32'$).

\subsection{Stellar Aberration}
The finite speed of light combined with the observer's own motion tilts the apparent direction of an incoming ray, exactly as vertically falling rain appears to slant forward when run through by a moving observer. Consider a photon arriving from directly overhead (velocity $-c\hat{\vec y}$ in the Sun's rest frame) while the Earth moves with velocity $v\hat{\vec x}$ transverse to that direction. Since $v \ll c$, Galilean velocity addition is an excellent approximation: in the Earth's frame, the photon's apparent velocity is
\begin{equation}
	\vec c\,' = \vec c - \vec v = -v\,\hat{\vec x} - c\,\hat{\vec y},
\end{equation}
so the ray appears tilted from the vertical by an angle $\kappa$ with
\begin{equation}
	\tan\kappa = \frac{v}{c} \quad\Longrightarrow\quad \kappa \approx \frac{v}{c} \quad (\text{since } v\ll c).
\end{equation}
Using Earth's orbital speed $v \approx 29.8\,\mathrm{km/s}$ (transverse to the line of sight to a star near the ecliptic pole, the maximal case) and converting to arcseconds:
\begin{keyresult}
	\begin{equation}
		\kappa = \frac{v}{c} \approx 20.4955'' \approx 20.5''.
		\label{eq:aberration_constant}
	\end{equation}
\end{keyresult}
Because $v$ is Earth's \emph{orbital} velocity, $\kappa$ varies sinusoidally over the year and traces out a small ellipse (or, for stars at the ecliptic pole, a circle) of angular semi-major axis $20.5''$ on the sky --- geometrically a much smaller, much faster analogue of the parallax ellipse of Chapter~3, and one that (unlike parallax) is completely independent of the star's distance.

\section{Angular Distance on a Sphere}

The angular separation $\theta$ between two points $(\alpha_1, \delta_1)$ and $(\alpha_2, \delta_2)$ is given by the dot product of their unit vectors (Eq.~\eqref{eq:unit_vector_eq}), since for any two unit vectors $\cos\theta = \hat{\mathbf r}_1\cdot\hat{\mathbf r}_2$:
\begin{align}
	\cos\theta &= (\cos\delta_1\cos\alpha_1)(\cos\delta_2\cos\alpha_2) + (\cos\delta_1\sin\alpha_1)(\cos\delta_2\sin\alpha_2) + \sin\delta_1\sin\delta_2 \notag \\
	&= \cos\delta_1\cos\delta_2\big(\cos\alpha_1\cos\alpha_2 + \sin\alpha_1\sin\alpha_2\big) + \sin\delta_1\sin\delta_2,
\end{align}
which, using $\cos\alpha_1\cos\alpha_2+\sin\alpha_1\sin\alpha_2=\cos(\alpha_1-\alpha_2)$, gives the spherical law of cosines:
\begin{keyresult}
	\begin{equation}
		\cos\theta = \hat{\mathbf r}_1 \cdot \hat{\mathbf r}_2 = \sin\delta_1\sin\delta_2 + \cos\delta_1\cos\delta_2\cos(\alpha_1 - \alpha_2).
		\label{eq:spherical_law_cos}
	\end{equation}
\end{keyresult}
For small angular separations ($\theta \ll 1$), the Haversine formula avoids floating-point precision loss:
\begin{equation}
	\sin^2\left(\frac{\theta}{2}\right) = \sin^2\left(\frac{\delta_2 - \delta_1}{2}\right) + \cos\delta_1\cos\delta_2\sin^2\left(\frac{\alpha_2 - \alpha_1}{2}\right).
	\label{eq:haversine_formula}
\end{equation}

\section{The Ecliptic Coordinate System}

\begin{definition}[Ecliptic Coordinates]
	\emph{Ecliptic longitude} $\lambda \in [0^\circ, 360^\circ)$ is measured eastward along the ecliptic from $\hat\gamma$. \emph{Ecliptic latitude} $\beta \in [-90^\circ, +90^\circ]$ is measured perpendicular to the ecliptic plane.
\end{definition}

\begin{figure}[htbp]
	\centering
	\includegraphics[width=0.4\linewidth]{Equator-Elliptic.png}
	\caption{Obliquity of the ecliptic $\epsilon \approx 23.44^\circ$ relating equatorial and ecliptic planes.}
	\label{fig:equator_elliptic}
\end{figure}

The ecliptic and equatorial planes intersect exactly along the line to the vernal equinox $\hat\gamma$ (by definition of $\hat\gamma$, Section~\ref{sec:sun_coordinates} of Chapter~1), so a rotation by the obliquity $\epsilon$ about that shared axis (the Cartesian $x$-axis in both frames) carries one frame into the other --- precisely the same kind of orthogonal-rotation-matrix argument used to derive the horizontal transformation above, only with a fixed rotation axis and angle rather than one set by the observer's latitude:
\begin{equation}
	\begin{pmatrix} \cos\beta\cos\lambda \\ \cos\beta\sin\lambda \\ \sin\beta \end{pmatrix} = \begin{pmatrix} 1 & 0 & 0 \\ 0 & \cos\epsilon & \sin\epsilon \\ 0 & -\sin\epsilon & \cos\epsilon \end{pmatrix} \begin{pmatrix} \cos\delta\cos\alpha \\ \cos\delta\sin\alpha \\ \sin\delta \end{pmatrix}.
	\label{eq:ecliptic_rotation_matrix}
\end{equation}
This yields:
\begin{keyresult}
	\begin{align}
		\sin\beta &= \sin\delta\cos\epsilon - \cos\delta\sin\alpha\sin\epsilon, \label{eq:ecliptic_beta} \\
		\cos\beta\cos\lambda &= \cos\delta\cos\alpha, \label{eq:ecliptic_coslam} \\
		\cos\beta\sin\lambda &= \cos\delta\sin\alpha\cos\epsilon + \sin\delta\sin\epsilon. \label{eq:ecliptic_sinlam}
	\end{align}
\end{keyresult}

\section{The Galactic Coordinate System}

For studies of Galactic structure and stellar kinematics, coordinates are defined relative to the disk of the Milky Way.

\begin{definition}[Galactic Coordinates $(l, b)$]
	\emph{Galactic longitude} $l \in [0^\circ, 360^\circ)$ is measured eastward along the Galactic equator from the Galactic Center (Sagittarius A*). \emph{Galactic latitude} $b \in [-90^\circ, +90^\circ]$ is the angle north ($+$) or south ($-$) of the Galactic midplane.
\end{definition}

By IAU definition (epoch J2000.0), the North Galactic Pole ($\mathrm{NGP}$) and Galactic Center are fixed at:
\begin{align}
	\text{North Galactic Pole (NGP):}\quad &\alpha_{\rm GP} = 12^{\rm h}51^{\rm m}26.28^{\rm s} = 192.8595^\circ, \quad \delta_{\rm GP} = +27^\circ 07' 42.0'' = +27.1283^\circ, \\
	\text{Position Angle of NCP at NGP:}\quad &\theta_0 = 122.932^\circ.
\end{align}

By exactly the same rotation-matrix method as above --- now rotating (and shifting the longitude origin) to align the polar axis with the Galactic pole $(\alpha_{\rm GP},\delta_{\rm GP})$ instead of the ecliptic pole --- the transformation from equatorial coordinates $(\alpha, \delta)$ to Galactic coordinates $(l, b)$ is:
\begin{keyresult}
	\begin{align}
		\sin b &= \sin\delta_{\rm GP}\sin\delta + \cos\delta_{\rm GP}\cos\delta\cos(\alpha - \alpha_{\rm GP}), \label{eq:gal_b} \\
		\cos b\sin(l_{\rm CP} - l) &= \cos\delta\sin(\alpha - \alpha_{\rm GP}), \label{eq:gal_sinl} \\
		\cos b\cos(l_{\rm CP} - l) &= \cos\delta_{\rm GP}\sin\delta - \sin\delta_{\rm GP}\cos\delta\cos(\alpha - \alpha_{\rm GP}), \label{eq:gal_cosl}
	\end{align}
\end{keyresult}
where $l_{\rm CP} = 122.932^\circ$ is the Galactic longitude of the North Celestial Pole.

\section{Airmass and the Parallactic Angle}

\subsection{Airmass ($X$)}
The \emph{airmass} $X$ measures the optical path length through the atmosphere relative to the zenith path. Model the atmosphere as a stack of thin, flat, horizontal slabs (the \emph{plane-parallel approximation}, valid away from the horizon where Earth's curvature is negligible). A ray arriving at zenith angle $z$ crosses a slab of vertical thickness $dz'$ along a slanted path of length $dz'/\cos z$ (elementary right-triangle geometry: the slant path is the hypotenuse of a triangle with vertical leg $dz'$ and angle $z$ from the vertical). Summing over all slabs,
\begin{equation}
	\int \rho(s)\,ds = \frac{1}{\cos z}\int \rho(z')\,dz',
\end{equation}
so the ratio of the slant path to the vertical (zenith) path is
\begin{keyresult}
	\begin{equation}
		X \equiv \frac{\int \rho(s)\,ds}{\int \rho(z)\,dz} \approx \sec z = \frac{1}{\cos z} = \frac{1}{\sin h}.
		\label{eq:airmass_secz}
	\end{equation}
\end{keyresult}
This plane-parallel result diverges unphysically as $z\to 90^\circ$; in reality Earth's curvature caps the path length at a finite value. For $z > 60^\circ$, this curvature becomes significant, and no simple geometric argument like the one above applies — instead one uses the empirical fitting formula of Kasten \& Young (1989):
\begin{equation}
	X = \frac{1}{\cos z + 0.50572\,(96.07995^\circ - z)^{-1.6364}}.
	\label{eq:airmass_kasten}
\end{equation}

\subsection{The Parallactic Angle ($q$)}
The \emph{parallactic angle} $q$ is the angle, measured at the star itself, between the great circle joining the star to the celestial pole and the great circle joining the star to the zenith. These two great circles, together with the great circle joining the pole to the zenith (an arc of length $90^\circ-\phi_0$), form a spherical triangle --- the \emph{astronomical triangle} --- with the star, the pole, and the zenith as its three vertices, and $H$ and $z=90^\circ-h$ as two of its known sides/angles. Applying the spherical laws of sines and cosines to this triangle (the same kind of spherical trigonometry that produced Eqs.~\eqref{eq:sinh_master}--\eqref{eq:cosh_cosA} for the pole-zenith-star triangle's other angle) yields $q$ directly:
\begin{equation}
	\sin q = \frac{\sin H\cos\phi_0}{\cos h}, \qquad \cos q = \frac{\sin\phi_0 - \sin h\sin\delta}{\cos h\cos\delta}.
	\label{eq:parallactic_angle}
\end{equation}
Aligning a spectrograph slit along the parallactic angle ($q$) minimizes atmospheric differential dispersion light loss across optical wavelengths.

Chapters~1 and~2 together fix \emph{when} and \emph{where} on the sky an observation is made. Chapter~3 now asks how far away the observed object actually is, and how it is moving -- the coordinates $(\alpha,\delta)$ derived here become, once combined with a distance, the angular half of a full three-dimensional position and velocity vector.
